\section{Conclusion}
\label{sec:conclusion}

QuaRK organizes frozen quantum reservoir learning around a compact hierarchy:
an architecture $\arch$ specifies the representation, a design
$\design\sim\designLaw_{\arch}$ fixes one random realization, the causal map
$\featureMap_{\design,\infty}$ exposes topology-local observables, and a
dimension-normalized, norm-controlled Mat\'ern readout performs regression.
The RMS feature metric keeps one kernel specification meaningful across
architectures without per-$(n,R)$ retuning.  A finite pool of such designs is
selected by full-sample ERM, with no optimization of internal quantum
parameters.

The theory first clarifies why the declared mixer is not merely decorative.
With product input injection, the identity mixer leaves every raw observable
coordinate confined to its seed qubits; generic graph-local mixing activates
higher-support Pauli components, and edge observables expose the first
nontrivial correlation layer and seed further spreading.  This is a
finite-resource feature statement rather than a claim of monotone benefit from
entanglement, and the empirical mixer ablation is designed to test whether the
additional observable structure is task useful.

The representation and learning results then connect directly.  At width $n$,
the projection-information floor isolates what the shared
Gaussian preprocessing can discard; it is nonincreasing under coupled widths
and vanishes almost surely in the injective regime.  QuaRK's finite-multiplex
support family can approach that floor, and continuity turns a strictly good
support design into positive random-design mass.  Pool size $S$ therefore has
an explicit probabilistic role through the miss probability
$(1-\hitMass)^S$, while the uniform beta-mixing bound controls the
overfitting price of selecting from that same pool.  Adding finite-shot and
finite-window fidelity yields a single excess-risk certificate for the
operational predictor returned by the algorithm.  For geometrically mixing
processes, the accompanying gap corollary makes the dependence requirement
explicit through a sufficient $O(\log N)$ raw-time gap.

A complementary finite-data JL theorem remains available for calibrating
observed preprocessing geometry, but it is no longer conflated with predictive
approximation.  The empirical program can therefore test the theory's intended
resource roles separately: when additional width preserves useful input
information, when the mixer and observable bank expose useful higher-order interactions,
when spatial multiplexing and random search improve frozen representations,
and when measurement or memory fidelity becomes the dominant cost.  The
reported acquisition metric counts branch-circuit runs under reusable $n$-qubit
hardware; per-run width and depth remain explicit separately through $n$ and
$w$.
